\centerline{\bf \Huge Ложка комбы в бочке ТЧ}

\begin{flushright}
        \scriptsize{Не знаю как но я запихнул алгебру, комбу, ТЧ и геому в одну домашку\\ и даже не назвал это разнобоем...}
\end{flushright}

\problem{1}
Найдите все такие пары натуральных чисел $m$ и $n$, что $m^{2025}+n$ делится на $m n$.

\problem{2}
Какое наибольшее количество натуральных чисел, не превосходящих 2n, можно
выбрать так, чтобы ни одно из них не делилось на другое?


\problem{3}
Даны натуральные числа $m$ и $n$ и два набора попарно различных вещественных чисел $a_1, \ldots, a_m$ и $b_1, \ldots, b_n$. Какое наименьшее количество различных чисел может встретиться среди $m n$ всевозможных сумм $a_i+b_j$ ?

\problem{4}
На доске написаны числа 2, 3, 5, ..., 2003, 2011, 2017, т.е. все простые числа, не превосходящие 2020. За одну операцию можно заменить два числа $a, b$ на максимальное простое число, не превосходящее $\sqrt{a^2-a b+b^2}$. После нескольких операций на доске осталось одно число. Какое максимальное значение оно может принимать?


\problem{5}\textbf{Двумерная теорема Минковского о выпуклом теле.} Начало координат является центром симметрии выпуклой фигуры площадью более 4. Докажите, что эта фигура содержит хотя бы одну точку с целыми координатами, отличную от начала координат.

\begin{flushright}
        \vspace{-5mm}
        \scriptsize{чел, какая ещё комбитчгеома???}
\end{flushright}

\problem{6}
Обозначим за $\sigma(n)$ сумму всех делителей числа $n$ (включая само число). Для каких $n$ выполняется неравенство $\sigma(8 n)>\sigma(9 n)$ ?